% Generated draft: requires mathematical and source review.
\documentclass[11pt]{article}
\usepackage[margin=25mm]{geometry}
\usepackage{fontspec}
\setmainfont{Latin Modern Roman}
\usepackage{amsmath,amssymb,mathtools,unicode-math}
\setmathfont{Latin Modern Math}
\usepackage{enumitem,needspace,xurl,hyperref}
\setlength{\parindent}{0pt}
\setlength{\parskip}{5pt}
\setlength{\emergencystretch}{3em}
\newcommand{\N}{\mathbb{N}}
\newcommand{\Z}{\mathbb{Z}}
\newcommand{\Q}{\mathbb{Q}}
\newcommand{\R}{\mathbb{R}}
\newcommand{\C}{\mathbb{C}}
\newcommand{\F}{\mathbb{F}}
\newcommand{\A}{\mathbb{A}}
\newcommand{\PP}{\mathbb{P}}
\newcommand{\E}{\mathbb{E}}
\newcommand{\eps}{\varepsilon}
\newcommand{\dd}{\,\mathrm{d}}
\begin{document}
\section*{UnsolvedMath ID: 10400106}
\textbf{Source problem code:} catalogue entry
\subsection{Definitions and mathematical statement}
A classical quandle is a set \(X\) with a binary operation \(*:X\times X\to X\) such that, for all \(x,y,z\in X\),
\[x*x=x,\qquad R_y:x\mapsto x*y\text{ is bijective},\qquad (x*y)*z=(x*z)*(y*z).\]
The problem attributed to M. Polyak is to define a \emph{quantum quandle}. This is a request to propose and justify an appropriate mathematical notion, not a conjecture about a structure already defined in the supplied record. A complete answer would have to state the underlying objects, operations and axioms and explain their relation to classical quandles. The source does not specify whether the desired relation is a deformation with a classical specialization, a noncommutative analogue, a categorical construction, or another quantization procedure. It supplies no coefficient ring, deformation parameter, classical-limit condition, morphisms, examples, or invariant-building requirement. Consequently the classical axioms above fix the object to be generalized but do not prescribe quantum axioms or assert that a particular existing meaning of the phrase meets the source's intention.
\subsection{Short English statement}
Polyak asks for a suitable definition of a quantum analogue of a quandle. The record does not specify what form of quantization or classical recovery the new definition should satisfy.
\subsection{Sources}
\begin{itemize}
\item Local UnsolvedMath registry, permanent ID \(10400106\), exact source code \texttt{catalogue entry}.
\item M. Polyak, attribution in Ohtsuki, \emph{Problems on invariants of knots and 3-manifolds} (2002), Problem 5.10, PDF page 92: \url{https://msp.org/gtm/2002/04/gtm-2002-04-024s.pdf}. Registry access 2026-07-29, extraction pdf-font-structure.
\end{itemize}
\end{document}
